% GENERATED FILE. Edit canonical lessons, then run npm run generate:books.
% canonical-source-hash: fnv1a64:332b06dc7ff236a2
% canonical-lessons: PT-C334-escorrer, PT-C334-tapar, PT-C334-dourar, PT-C334-gratinar, PT-C334-rechear, PT-R334-revisao-333-334, PT-R334-revisao-final-334

\chapter{Até ficar dourado}
\label{ch:pt-ate-ficar-dourado}
\hlchaptermodality{334}

\begin{chapteropening}
\textbf{By the end of this chapter:} \emph{I can describe the last steps of a recipe: draining, covering, browning, grilling the top and stuffing.}

Complete the last lesson of chapter 334: I can describe the last steps of a recipe: draining, covering, browning, grilling the top and stuffing.
\end{chapteropening}

\section[escorrer]{escorrer — to drain}

\label{lesson:PT-C334-escorrer}



\begin{quote}
Before the new one: say the Portuguese for to grease, then the Portuguese for to sprinkle.
\end{quote}

\subsection*{You'll want to know: escorrer}
\textbf{escorrer} — \textquotedblleft{}to drain\textquotedblright{}.

\begin{quote}
Ela escorreu a massa e juntou o molho.
\end{quote}

\begin{quote}
\emph{She drained the pasta and added the sauce.}
\end{quote}

\emph{Pôr a escorrer} — to leave to drain.

\begin{cousinweb}
\emph{escorrer} comes from Latin \emph{excurrere} \textquotedblleft{}to run out\textquotedblright{} (EtymDB), from \emph{ex} \textquotedblleft{}out\textquotedblright{} and \emph{currere} \textquotedblleft{}to run\textquotedblright{}.
\end{cousinweb}

\begin{grammarlens}[title={What you've built}]
One more word for telling, explaining and describing at B1.
\end{grammarlens}

\subsection*{Your turn}
\begin{itemize}
\raggedright
  \item \emph{Say it:} \emph{escorrer}
  \item \emph{Say it:} \emph{escorrer}, once more
  \item \emph{Say it:} a sentence of your own with \emph{escorrer}
  \item \emph{Recall:} say the Portuguese for to grease, then the Portuguese for to sprinkle, then say \emph{escorrer} again
\end{itemize}

\begin{tcolorbox}[breakable,colback=teal!4,colframe=teal!35!black,title={Before you move on}]
What does escorrer mean? (\textquotedblleft{}To drain\textquotedblright{}.) Say it once more.
\end{tcolorbox}
\section[tapar]{tapar — to cover (with a lid)}

\label{lesson:PT-C334-tapar}



\begin{quote}
Before the new one: say the Portuguese for to sprinkle, then the Portuguese for to drain.
\end{quote}

\subsection*{You'll want to know: tapar}
\textbf{tapar} — \textquotedblleft{}to cover (with a lid)\textquotedblright{}.

\begin{quote}
O cozinheiro tapou o tacho e esperou.
\end{quote}

\begin{quote}
\emph{The cook put the lid on the pot and waited.}
\end{quote}

Also to cover anything: \emph{tapar os olhos}.

\begin{cousinweb}
\emph{tapar} is usually traced to a Germanic word for a stopper or plug, like English \emph{tap}; the origin is not settled, and EtymDB has no entry.
\end{cousinweb}

\begin{grammarlens}[title={What you've built}]
One more word for telling, explaining and describing at B1.
\end{grammarlens}

\subsection*{Your turn}
\begin{itemize}
\raggedright
  \item \emph{Say it:} \emph{tapar}
  \item \emph{Say it:} \emph{tapar}, once more
  \item \emph{Say it:} a sentence of your own with \emph{tapar}
  \item \emph{Recall:} say the Portuguese for to sprinkle, then the Portuguese for to drain, then say \emph{tapar} again
\end{itemize}

\begin{tcolorbox}[breakable,colback=teal!4,colframe=teal!35!black,title={Before you move on}]
What does tapar mean? (\textquotedblleft{}To cover (with a lid)\textquotedblright{}.) Say it once more.
\end{tcolorbox}
\section[dourar]{dourar — to brown, to turn golden; to gild}

\label{lesson:PT-C334-dourar}



\begin{quote}
Before the new one: say the Portuguese for to drain, then the Portuguese for to cover.
\end{quote}

\subsection*{You'll want to know: dourar}
\textbf{dourar} — \textquotedblleft{}to brown, to turn golden; to gild\textquotedblright{}.

\begin{quote}
A cebola dourou em cinco minutos.
\end{quote}

\begin{quote}
\emph{The onion turned golden in five minutes.}
\end{quote}

\emph{Até dourar} — until golden.

\begin{cousinweb}
\emph{dourar} comes down from Old Portuguese (EtymDB), from Late Latin \emph{deaurāre} \textquotedblleft{}to gild\textquotedblright{}, built on \emph{aurum} \textquotedblleft{}gold\textquotedblright{}.
\end{cousinweb}

\begin{grammarlens}[title={What you've built}]
One more word for telling, explaining and describing at B1.
\end{grammarlens}

\subsection*{Your turn}
\begin{itemize}
\raggedright
  \item \emph{Say it:} \emph{dourar}
  \item \emph{Say it:} \emph{dourar}, once more
  \item \emph{Say it:} a sentence of your own with \emph{dourar}
  \item \emph{Recall:} say the Portuguese for to drain, then the Portuguese for to cover, then say \emph{dourar} again
\end{itemize}

\begin{tcolorbox}[breakable,colback=teal!4,colframe=teal!35!black,title={Before you move on}]
What does dourar mean? (\textquotedblleft{}To brown, to turn golden; to gild\textquotedblright{}.) Say it once more.
\end{tcolorbox}
\section[gratinar]{gratinar — to brown on top in the oven, to gratinate}

\label{lesson:PT-C334-gratinar}



\begin{quote}
Before the new one: say the Portuguese for to cover, then the Portuguese for to brown.
\end{quote}

\subsection*{You'll want to know: gratinar}
\textbf{gratinar} — \textquotedblleft{}to brown on top in the oven, to gratinate\textquotedblright{}.

\begin{quote}
No fim, gratinamos o bacalhau no forno.
\end{quote}

\begin{quote}
\emph{At the end we brown the salt cod in the oven.}
\end{quote}

With cheese or breadcrumbs on top.

\begin{cousinweb}
\emph{gratinar} was borrowed from French \emph{gratiner}, from \emph{gratin} \textquotedblleft{}the crust\textquotedblright{}, built on \emph{gratter} \textquotedblleft{}to scrape\textquotedblright{}, a verb of Germanic origin (EtymDB has no entry).
\end{cousinweb}

\begin{grammarlens}[title={What you've built}]
One more word for telling, explaining and describing at B1.
\end{grammarlens}

\subsection*{Your turn}
\begin{itemize}
\raggedright
  \item \emph{Say it:} \emph{gratinar}
  \item \emph{Say it:} \emph{gratinar}, once more
  \item \emph{Say it:} a sentence of your own with \emph{gratinar}
  \item \emph{Recall:} say the Portuguese for to cover, then the Portuguese for to brown, then say \emph{gratinar} again
\end{itemize}

\begin{tcolorbox}[breakable,colback=teal!4,colframe=teal!35!black,title={Before you move on}]
What does gratinar mean? (\textquotedblleft{}To brown on top in the oven, to gratinate\textquotedblright{}.) Say it once more.
\end{tcolorbox}
\section[rechear]{rechear — to stuff, to fill}

\label{lesson:PT-C334-rechear}



\begin{quote}
Before the new one: say the Portuguese for to brown, then the Portuguese for to brown on top in the oven.
\end{quote}

\subsection*{You'll want to know: rechear}
\textbf{rechear} — \textquotedblleft{}to stuff, to fill\textquotedblright{}.

\begin{quote}
A avó recheou o frango com pão e chouriço.
\end{quote}

\begin{quote}
\emph{Grandmother stuffed the chicken with bread and chouriço.}
\end{quote}

\emph{O recheio} — the stuffing, the filling.

\begin{cousinweb}
\emph{rechear} is built from \emph{re-} and \emph{cheio} \textquotedblleft{}full\textquotedblright{}, which is inherited from Latin \emph{plēnus} \textquotedblleft{}full\textquotedblright{} (Lewis \& Short); EtymDB has the word with no source.
\end{cousinweb}

\begin{grammarlens}[title={What you've built}]
One more word for telling, explaining and describing at B1.
\end{grammarlens}

\subsection*{Your turn}
\begin{itemize}
\raggedright
  \item \emph{Say it:} \emph{rechear}
  \item \emph{Say it:} \emph{rechear}, once more
  \item \emph{Say it:} a sentence of your own with \emph{rechear}
  \item \emph{Recall:} say the Portuguese for to brown, then the Portuguese for to brown on top in the oven, then say \emph{rechear} again
\end{itemize}

\begin{tcolorbox}[breakable,colback=teal!4,colframe=teal!35!black,title={Before you move on}]
What does rechear mean? (\textquotedblleft{}To stuff, to fill\textquotedblright{}.) Say it once more.
\end{tcolorbox}
\section[(dialogue)]{Revisão: passo a passo}

\label{lesson:PT-R334-revisao-333-334}



\begin{quote}
Say the Portuguese for the electric mixer and to sift.
\end{quote}

\subsection*{Your turn}
Say these aloud:

\begin{itemize}
\raggedright
  \item the electric mixer, to sift, to knead, to grease, to sprinkle
  \item to drain, to cover, to brown, to brown on top in the oven, to stuff
\end{itemize}

\begin{tcolorbox}[breakable,colback=teal!4,colframe=teal!35!black,title={Before you move on}]
The electric mixer? (\textbf{a batedeira}.) To sift? (\textbf{peneirar}.) To knead? (\textbf{amassar}.) To grease? (\textbf{untar}.) To sprinkle? (\textbf{polvilhar}.) To drain? (\textbf{escorrer}.) To cover? (\textbf{tapar}.) To brown? (\textbf{dourar}.) To brown on top in the oven? (\textbf{gratinar}.) To stuff? (\textbf{rechear}.)
\end{tcolorbox}
\section[(dialogue)]{Até ficar dourado — once more}

\label{lesson:PT-R334-revisao-final-334}



\begin{quote}
Say the Portuguese for to drain and to cover.
\end{quote}

\subsection*{Your turn}
\begin{itemize}
\raggedright
  \item \emph{Say it:} to drain, to cover, to brown, to brown on top in the oven, to stuff
\end{itemize}

\begin{tcolorbox}[breakable,colback=teal!4,colframe=teal!35!black,title={Before you move on}]
To drain? (\textbf{escorrer}.) To cover? (\textbf{tapar}.) To brown? (\textbf{dourar}.) To brown on top in the oven? (\textbf{gratinar}.) To stuff? (\textbf{rechear}.)
\end{tcolorbox}
